% SpqrBridge slides, alternatives.
% Build: latexmk -lualatex slides.tex
% Shared preamble of the SpqrBridge slide decks.
% Same style and colors as ref-impl-simulation/slides/methodology.
% Each deck loads it with % Shared preamble of the SpqrBridge slide decks.
% Same style and colors as ref-impl-simulation/slides/methodology.
% Each deck loads it with % Shared preamble of the SpqrBridge slide decks.
% Same style and colors as ref-impl-simulation/slides/methodology.
% Each deck loads it with \input{../common/preamble} and builds with: latexmk -lualatex slides.tex
\documentclass[aspectratio=169,11pt]{beamer}

\usepackage{fontspec}
\setsansfont{DejaVu Sans}
\directlua{
  luaotfload.add_fallback("leanfallback", {"NotoSansMath:mode=harf;", "STIX Two Math:mode=harf;"})
}
\setmonofont{DejaVu Sans Mono}[Scale=0.9, RawFeature={fallback=leanfallback}]

\usepackage{tikz}
\usepackage{array}
\usetikzlibrary{arrows.meta, positioning, calc, fit, backgrounds, shapes.misc, decorations.pathreplacing}

% Colors. Spec side, impl side, proof, trusted, gap.
\definecolor{refc}{HTML}{2B6CB0}
\definecolor{realc}{HTML}{C05621}
\definecolor{proofc}{HTML}{2F855A}
\definecolor{trustc}{HTML}{718096}
\definecolor{gapc}{HTML}{C53030}
\definecolor{ink}{HTML}{1A202C}
\definecolor{paper}{HTML}{FFFFFF}

\setbeamercolor{normal text}{fg=ink, bg=paper}
\setbeamercolor{frametitle}{fg=ink}
\setbeamercolor{title}{fg=ink}
\setbeamercolor{subtitle}{fg=trustc}
\setbeamercolor{structure}{fg=refc}
\setbeamerfont{frametitle}{size=\Large, series=\bfseries}
\setbeamerfont{subtitle}{size=\footnotesize}
\setbeamertemplate{navigation symbols}{}
\setbeamertemplate{footline}{%
  \hfill\textcolor{trustc}{\scriptsize\insertframenumber/\inserttotalframenumber}\hspace{1.2em}\vspace{0.8em}}

\tikzset{
  >={Stealth[length=2.4mm]},
  box/.style={draw, rounded corners=3pt, thick, align=center, inner sep=5pt, font=\small},
  spec/.style={box, draw=refc, fill=refc!8},
  impl/.style={box, draw=realc, fill=realc!8},
  proof/.style={box, draw=proofc, fill=proofc!10},
  trust/.style={box, draw=trustc, fill=trustc!10},
  gap/.style={box, draw=gapc, fill=gapc!6, dashed},
  tool/.style={box, draw=ink!60, fill=ink!4},
  doc/.style={box, draw=ink!60, fill=paper, double, double distance=1pt},
  flow/.style={->, thick, ink!70},
  proved/.style={<->, very thick, proofc},
  rel/.style={-, very thick, proofc, dashed},
  label/.style={font=\scriptsize, align=center},
}

\newcommand{\code}[1]{\texttt{#1}}
\newcommand{\specside}[1]{\textcolor{refc}{#1}}
\newcommand{\implside}[1]{\textcolor{realc}{#1}}

\author{SpqrBridge}
\date{November 2026}
 and builds with: latexmk -lualatex slides.tex
\documentclass[aspectratio=169,11pt]{beamer}

\usepackage{fontspec}
\setsansfont{DejaVu Sans}
\directlua{
  luaotfload.add_fallback("leanfallback", {"NotoSansMath:mode=harf;", "STIX Two Math:mode=harf;"})
}
\setmonofont{DejaVu Sans Mono}[Scale=0.9, RawFeature={fallback=leanfallback}]

\usepackage{tikz}
\usepackage{array}
\usetikzlibrary{arrows.meta, positioning, calc, fit, backgrounds, shapes.misc, decorations.pathreplacing}

% Colors. Spec side, impl side, proof, trusted, gap.
\definecolor{refc}{HTML}{2B6CB0}
\definecolor{realc}{HTML}{C05621}
\definecolor{proofc}{HTML}{2F855A}
\definecolor{trustc}{HTML}{718096}
\definecolor{gapc}{HTML}{C53030}
\definecolor{ink}{HTML}{1A202C}
\definecolor{paper}{HTML}{FFFFFF}

\setbeamercolor{normal text}{fg=ink, bg=paper}
\setbeamercolor{frametitle}{fg=ink}
\setbeamercolor{title}{fg=ink}
\setbeamercolor{subtitle}{fg=trustc}
\setbeamercolor{structure}{fg=refc}
\setbeamerfont{frametitle}{size=\Large, series=\bfseries}
\setbeamerfont{subtitle}{size=\footnotesize}
\setbeamertemplate{navigation symbols}{}
\setbeamertemplate{footline}{%
  \hfill\textcolor{trustc}{\scriptsize\insertframenumber/\inserttotalframenumber}\hspace{1.2em}\vspace{0.8em}}

\tikzset{
  >={Stealth[length=2.4mm]},
  box/.style={draw, rounded corners=3pt, thick, align=center, inner sep=5pt, font=\small},
  spec/.style={box, draw=refc, fill=refc!8},
  impl/.style={box, draw=realc, fill=realc!8},
  proof/.style={box, draw=proofc, fill=proofc!10},
  trust/.style={box, draw=trustc, fill=trustc!10},
  gap/.style={box, draw=gapc, fill=gapc!6, dashed},
  tool/.style={box, draw=ink!60, fill=ink!4},
  doc/.style={box, draw=ink!60, fill=paper, double, double distance=1pt},
  flow/.style={->, thick, ink!70},
  proved/.style={<->, very thick, proofc},
  rel/.style={-, very thick, proofc, dashed},
  label/.style={font=\scriptsize, align=center},
}

\newcommand{\code}[1]{\texttt{#1}}
\newcommand{\specside}[1]{\textcolor{refc}{#1}}
\newcommand{\implside}[1]{\textcolor{realc}{#1}}

\author{SpqrBridge}
\date{November 2026}
 and builds with: latexmk -lualatex slides.tex
\documentclass[aspectratio=169,11pt]{beamer}

\usepackage{fontspec}
\setsansfont{DejaVu Sans}
\directlua{
  luaotfload.add_fallback("leanfallback", {"NotoSansMath:mode=harf;", "STIX Two Math:mode=harf;"})
}
\setmonofont{DejaVu Sans Mono}[Scale=0.9, RawFeature={fallback=leanfallback}]

\usepackage{tikz}
\usepackage{array}
\usetikzlibrary{arrows.meta, positioning, calc, fit, backgrounds, shapes.misc, decorations.pathreplacing}

% Colors. Spec side, impl side, proof, trusted, gap.
\definecolor{refc}{HTML}{2B6CB0}
\definecolor{realc}{HTML}{C05621}
\definecolor{proofc}{HTML}{2F855A}
\definecolor{trustc}{HTML}{718096}
\definecolor{gapc}{HTML}{C53030}
\definecolor{ink}{HTML}{1A202C}
\definecolor{paper}{HTML}{FFFFFF}

\setbeamercolor{normal text}{fg=ink, bg=paper}
\setbeamercolor{frametitle}{fg=ink}
\setbeamercolor{title}{fg=ink}
\setbeamercolor{subtitle}{fg=trustc}
\setbeamercolor{structure}{fg=refc}
\setbeamerfont{frametitle}{size=\Large, series=\bfseries}
\setbeamerfont{subtitle}{size=\footnotesize}
\setbeamertemplate{navigation symbols}{}
\setbeamertemplate{footline}{%
  \hfill\textcolor{trustc}{\scriptsize\insertframenumber/\inserttotalframenumber}\hspace{1.2em}\vspace{0.8em}}

\tikzset{
  >={Stealth[length=2.4mm]},
  box/.style={draw, rounded corners=3pt, thick, align=center, inner sep=5pt, font=\small},
  spec/.style={box, draw=refc, fill=refc!8},
  impl/.style={box, draw=realc, fill=realc!8},
  proof/.style={box, draw=proofc, fill=proofc!10},
  trust/.style={box, draw=trustc, fill=trustc!10},
  gap/.style={box, draw=gapc, fill=gapc!6, dashed},
  tool/.style={box, draw=ink!60, fill=ink!4},
  doc/.style={box, draw=ink!60, fill=paper, double, double distance=1pt},
  flow/.style={->, thick, ink!70},
  proved/.style={<->, very thick, proofc},
  rel/.style={-, very thick, proofc, dashed},
  label/.style={font=\scriptsize, align=center},
}

\newcommand{\code}[1]{\texttt{#1}}
\newcommand{\specside}[1]{\textcolor{refc}{#1}}
\newcommand{\implside}[1]{\textcolor{realc}{#1}}

\author{SpqrBridge}
\date{November 2026}


\title{Is there a simpler approach?}
\subtitle{SpqrBridge slides, alternatives}

\begin{document}

\begin{frame}
  \titlepage
  \centering
  {\footnotesize\textcolor{trustc}{Can we map \code{send}, \code{recv} and \code{init} function by function, with no state machine?}}
\end{frame}

\begin{frame}[t]{Can we map \code{send} to \code{send}, and stop there?}\label{fr:square}
  \footnotesize
  \code{send\_chunked} is that mapping, from the extracted \code{send} to
  \code{SPQR.Chunked.send}. Its statement needs three facts about states.

  \centering
  \begin{tikzpicture}[scale=0.8, transform shape, every node/.append style={font=\scriptsize}]
    % one square
    \node[impl, minimum width=1.5cm] (s) at (0, 1.0) {$s$};
    \node[impl, minimum width=1.5cm] (s') at (4.2, 1.0) {$s',\ m,\ k$};
    \node[spec, minimum width=1.5cm] (t) at (0, -0.6) {$t$};
    \node[spec, minimum width=1.5cm] (t') at (4.2, -0.6) {$t',\ m',\ k'$};
    \draw[flow, realc] (s) -- node[above, label, text=realc] {code \code{send}} (s');
    \draw[flow, refc] (t) -- node[below, label, text=refc] {\code{Chunked.send}} (t');
    \draw[rel] (s) -- node[left, label, text=proofc] {$RA$} (t);
    \draw[rel] (s') -- node[right, label, text=proofc] {$RA$, $RM$, $=$} (t');

    \node[gap, font=\scriptsize, text width=6.4cm, align=left, anchor=west] (p1) at (6.6, 1.2)
      {\textbf{1.} $RA\ s\ t$: related inputs. The types differ, so equality is not possible.};
    \node[gap, font=\scriptsize, text width=6.4cm, align=left, anchor=west] (p2) at (6.6, 0.25)
      {\textbf{2.} \code{ReachParty} $s$: a run from \code{init} ends in $s$.};
    \node[gap, font=\scriptsize, text width=6.4cm, align=left, anchor=west] (p3) at (6.6, -0.55)
      {\textbf{3.} The epoch of $s$ is below the \code{u64} limit.};
  \end{tikzpicture}

  \begin{tikzpicture}[scale=0.8, transform shape, every node/.append style={font=\scriptsize},
      c/.style={impl, minimum width=0.8cm, minimum height=0.5cm, inner sep=2pt},
      m/.style={spec, minimum width=0.8cm, minimum height=0.5cm, inner sep=2pt}]
    \foreach \i/\l in {0/{s_0}, 1/{s_1}, 2/{s_2}, 3/{s_3}, 4/{s_4}} {
      \node[c] (c\i) at (2.3*\i, 0.5) {$\l$};
    }
    \foreach \i/\l in {0/{t_0}, 1/{t_1}, 2/{t_2}, 3/{t_3}, 4/{t_4}} {
      \node[m] (m\i) at (2.3*\i, -0.5) {$\l$};
      \draw[rel] (c\i) -- (m\i);
    }
    \foreach \i/\j/\l in {0/1/send, 1/2/recv, 2/3/send, 3/4/recv} {
      \draw[flow, realc] (c\i) -- node[above, label, text=realc] {\l} (c\j);
      \draw[flow, refc] (m\i) -- node[below, label, text=refc] {\l} (m\j);
    }
    \node[label, text=trustc, anchor=east] at (-0.7, 0.5) {\code{init}};
    \node[label, text=trustc, anchor=east] at (-0.7, -0.5) {\code{init}};
    \node[label, text=proofc, anchor=west] at (9.9, 0) {$RA$ is the\\induction\\hypothesis};
  \end{tikzpicture}
  \vspace{0.2em}

  \raggedright
  \begin{itemize}\footnotesize
    \item The output of a send depends on its input state, and that state depends on the earlier
      calls. So the proof carries $RA$ from \code{init} to step $n$.
    \item One square for each function, chained along a run, is a simulation between two state
      machines. The per-function mapping is one step of it.
  \end{itemize}
\end{frame}

\begin{frame}[t]{The mapping holds on reachable states only}\label{fr:reach}
  \footnotesize
  For an arbitrary pair of related states, the square of slide \ref{fr:square} is false or not
  claimed. The premises that repair it are facts about runs.

  \centering
  \begin{tikzpicture}[scale=0.78, transform shape, every node/.append style={font=\scriptsize}]
    \draw[trustc, thick, rounded corners=8pt, fill=trustc!6] (-6.9, -2.3) rectangle (6.9, 2.3);
    \node[label, text=trustc, anchor=north west] at (-6.8, 2.25) {all values of the type \code{States}};
    \draw[proofc, very thick, fill=proofc!10] (0, -0.15) ellipse (2.7cm and 1.45cm);
    \node[label, text=proofc] at (0, 0.85) {\textbf{reachable}: \code{CodeReach}};
    \node[impl, inner sep=2pt, font=\tiny] (i) at (-1.7, -0.2) {\code{init}};
    \node[impl, inner sep=2pt, font=\tiny] (a) at (-0.55, -0.2) {$s_1$};
    \node[impl, inner sep=2pt, font=\tiny] (b) at (0.55, -0.2) {$s_2$};
    \node[impl, inner sep=2pt, font=\tiny] (c) at (1.7, -0.2) {$s_3$};
    \draw[flow, realc] (i) -- (a);
    \draw[flow, realc] (a) -- (b);
    \draw[flow, realc] (b) -- (c);
    \node[label, text=proofc] at (0, -0.95) {the squares hold here};

    \node[gap, font=\scriptsize, text width=3.5cm, align=left] (g1) at (-4.8, 0.85)
      {epoch $= 2^{64} - 1$:\\the code refuses,\\SM continues (main deck)};
    \node[gap, font=\scriptsize, text width=3.5cm, align=left] (g2) at (-4.8, -1.15)
      {a state that no run produces:\\no claim about \code{send}\\(\#785, \#787)};
    \node[gap, font=\scriptsize, text width=3.5cm, align=left] (g3) at (4.8, 0.85)
      {a message that the other\\party never sent:\\no claim about \code{recv}};
    \node[gap, font=\scriptsize, text width=3.5cm, align=left] (g4) at (4.8, -1.15)
      {SM refuses a message:\\the code can do anything};
  \end{tikzpicture}
  \vspace{0.2em}

  \raggedright
  \begin{itemize}\footnotesize
    \item \code{CodeReach} is inductive over \code{init}, \code{send} and \code{recv} of both
      parties. Without the machine, ``reachable'' has no meaning.
    \item The receive premise needs both parties: the message is in the list that the
      \emph{other} party sent in the same run (\code{ReachRecv}).
    \item The bottom-right case is reachable. The game permits it, because SM's game failed
      already. An equality of two functions cannot state this.
  \end{itemize}
\end{frame}

\begin{frame}[t]{What needs states, and what is a choice}
  \vspace{0.6em}
  \centering
  \scriptsize
  \renewcommand{\arraystretch}{1.3}
  \begin{tabular}{>{\raggedright\arraybackslash}p{4.6cm} p{2.1cm} >{\raggedright\arraybackslash}p{6.6cm}}
    \textbf{piece} & \textbf{needed?} & \textbf{why} \\
    \hline
    a relation between code states and SM states ($RA$)
      & \textcolor{gapc}{yes} & the types differ, and \code{send} takes a state and returns one
      (slide \ref{fr:square}) \\
    a reachability invariant (\code{CodeReach})
      & \textcolor{gapc}{yes} & the step facts fail or are not claimed outside it
      (slide \ref{fr:reach}) \\
    a game-level argument (invariant, potential)
      & \textcolor{gapc}{yes} & the property is about runs of two parties
      (run deck) \\
    a state map for corruptions (\code{conc})
      & \textcolor{gapc}{for security} & the adversary reads the state
      (run deck) \\
    \code{Braid} between the code and SM
      & \textcolor{proofc}{a choice} & one direct relation is possible (route A, next slides) \\
    the eleven named states
      & \textcolor{proofc}{not ours} & the Rust type \code{States} is an enum, and
      \code{send} matches on it \\
  \end{tabular}
  \vspace{0.3em}

  \raggedright\footnotesize
  \begin{itemize}
    \item The proposal ``map \code{send} to \code{send}, \code{recv} to \code{recv}, \code{init}
      to \code{init}'' is route A. It removes \code{Braid}. It keeps the first four rows.
    \item The bridge already contains that mapping: \code{send\_chunked}, \code{receive\_chunked}
      and the \code{init} obligation. The transfer theorem takes exactly these three as input.
  \end{itemize}
\end{frame}

\begin{frame}[t]{Two routes from SM's bound to the code}
  \footnotesize
  SM proves its bound for the game over \code{SPQR.Chunked}. To move it, relate each state of the
  code's game to a state of SM's game, step by step.
  \vspace{0.1em}

  \centering
  \begin{tikzpicture}[every node/.append style={font=\scriptsize}]
    \node[impl, minimum width=3.0cm] (c) at (0, 0) {\textbf{code}\\extracted \code{States}};
    \node[spec, minimum width=3.0cm] (k) at (9, 0) {\textbf{\code{SPQR.Chunked}}\\SM model};
    \node[spec, minimum width=3.0cm] (b) at (4.5, 2.1) {\textbf{\code{Braid}}\\ML-KEM Braid, with\\the code's deviations};
    \node[proof, below=0.7cm of k] (q) {$\Pr[\text{fail}] \le q\varepsilon$};
    \draw[flow, proofc] (k) -- node[right, label, text=proofc] {SM proof} (q);

    \draw[gapc, very thick, dashed] (c) -- node[below=4pt, label, text=gapc]
      {\textbf{route A}: one direct relation} (k);

    \draw[rel] (c) -- node[left=6pt, label, text=proofc, pos=0.5]
      {\code{Refines}\\spqr-verify} (b.west);
    \draw[rel] (b.east) -- node[right=6pt, label, text=proofc, pos=0.5]
      {\code{ChunkedRel}\\this repository} (k);
    \node[label, text=proofc] at (4.5, 0.75) {\textbf{route B}: through \code{Braid}};
  \end{tikzpicture}
  \vspace{0.3em}

  \raggedright
  \begin{itemize}\footnotesize
    \item Both routes must give one relation $RA$ between code states and SM states. Sends
      are coupled, and receives keep it until SM fails.
    \item Route B builds it as a composition through \code{Braid} (main deck).
  \end{itemize}
\end{frame}

\begin{frame}[t]{What a direct relation would have to prove}
  \footnotesize
  A step of route A relates one Rust step to one SM step. It mixes two kinds of facts.
  \vspace{0.4em}

  \centering
  \begin{tikzpicture}[every node/.append style={font=\scriptsize}]
    \node[gap, minimum width=12.4cm] (a) at (0, 0)
      {route A: \implside{one Rust step} $\sim$ \specside{one \code{SPQR.Chunked} step}};

    \node[impl, text width=6.2cm, align=left, font=\scriptsize] (l) at (-3.45, -0.6) [anchor=north] {%
      \textbf{about the code}
      \begin{itemize}\scriptsize\setlength{\itemsep}{0pt}
        \item \code{Result}: panics, divergence, \code{Err}
        \item \code{u64} and \code{usize} bounds
        \item bytes, slices, serialization
        \item calls to ML-KEM, HKDF, HMAC
        \item which random bytes each send reads
      \end{itemize}};
    \node[spec, text width=6.2cm, align=left, font=\scriptsize] (r) at (3.45, -0.6) [anchor=north] {%
      \textbf{about the models}
      \begin{itemize}\scriptsize\setlength{\itemsep}{0pt}
        \item byte fields against typed values
        \item epochs: \code{u64} against $\mathbb{N}$
        \item streaming decoder against set decoder
        \item encoder counter modulo $2^{32}$
        \item deviations of the code from §2.5
      \end{itemize}};
    \draw[flow, gapc] (a.south -| l) -- (l);
    \draw[flow, gapc] (a.south -| r) -- (r);

    \node[proof, text width=6.2cm, font=\scriptsize, below=0.3cm of l] (lp)
      {functional correctness: spqr-verify\\\code{Refines}, \#785, \#787};
    \node[proof, text width=6.2cm, font=\scriptsize, below=0.3cm of r] (rp)
      {model matching: this repository\\\code{ChunkedRel}};
    \draw[flow, proofc] (l) -- (lp);
    \draw[flow, proofc] (r) -- (rp);
  \end{tikzpicture}
  \vspace{0.4em}

  \raggedright
  Route A needs the left column too. It can reuse upstream lemmas, but it states each step
  against SM's types. Route B states these facts once, against \code{Braid}. Totality and the
  randomness fact are still premises there (\#785, \#787).
\end{frame}

\begin{frame}[t]{Why \code{Braid} in the middle}
  \centering
  \begin{tikzpicture}[every node/.append style={font=\scriptsize}]
    \node[spec, minimum width=2.8cm] (b) at (0, 0) {\textbf{\code{Braid}}\\one abstract\\machine};

    \node[spec, minimum width=3.0cm] (p) at (-5.0, 1.1) {ML-KEM Braid §2.5\\(the paper)};
    \node[impl, minimum width=3.0cm] (c) at (-5.0, -0.3) {code, v1.6.0};
    \node[impl, minimum width=3.0cm, dashed] (c2) at (-5.0, -1.25) {next code version};
    \node[spec, minimum width=3.0cm] (k) at (5.0, 1.1) {\code{SPQR.Chunked}\\SM correctness};
    \node[spec, minimum width=3.0cm, dashed] (s) at (5.0, -0.15) {other SM models\\or security games};
    \node[spec, minimum width=3.0cm, dashed] (o) at (5.0, -1.25) {other properties};

    \draw[ink!60, thick, dotted] (p.east) -- node[above, sloped, label, text=ink!70, pos=0.6] {deviations named} (b.north west);
    \draw[rel] (c.east) -- node[below, sloped, label, text=proofc, pos=0.55] {\code{Refines}} (b.west);
    \draw[rel, proofc!40] (c2.east) -- (b.south west);
    \draw[rel] (b.north east) -- node[above, sloped, label, text=proofc, pos=0.45] {\code{ChunkedRel}} (k.west);
    \draw[rel, proofc!40] (b.east) -- (s.west);
    \draw[rel, proofc!40] (b.south east) -- (o.west);

    \node[label, text=realc] at (-5.0, -1.85) {code side: spqr-verify};
    \node[label, text=refc] at (5.0, -1.85) {model side: bridges};
  \end{tikzpicture}
  \vspace{0.3em}

  \raggedright
  \begin{itemize}\footnotesize
    \item Each repository proves what it owns: facts about the Rust code next to the extraction,
      facts about SM's model next to SM.
    \item If \code{Braid}'s interface stays stable, a change on one side redoes one relation only:
      a new code version redoes \code{Refines}, a new model needs a relation from \code{Braid}.
    \item Security games are future work. They see leaked randomness, so the coupling must also relate
      the bytes that each side leaks.
    \item \code{Braid} reads like the paper, with each deviation named.
    \item The cost: one more model to maintain. \code{ChunkedRel} still needs two code facts: the
      Rust decoder is lawful, and the encoder counter wraps.
  \end{itemize}
\end{frame}

\end{document}
